% STATUS: written 2026-11-05, recalled results as is (one environment each, source kind and whole statement)
% Sources read in full: Spiegel 1994 (pp. 2955--2959), Ferrero, Giambruno and Polcino Milies 1995 (pp. 434--437).
% Additions of 2026-11-07 moved from chapter 20: Fox 1953 (pp. 549--550) read in full.
% Additions of 2026-11-06 from Arutyunov--Zhiltsov 2023 (pp. 1, 8),
% Arutyunov--Naianzin 2023 (pp. 1--2) and Antrobus--Briggs--Rubio y Degrassi 2026 (pp. 1--3).
% Sources: Spiegel1994, FerreroGiambrunoPolcinoMilies1995, ManjuSharma2026,
% ArutyunovZhiltsov2023, ArutyunovNaianzin2023, AntrobusBriggsRubioDegrassi2026.
\section{Group algebras and group rings}\label{sec:group-hopf:group-algebras}

Group rings join a group to a ring, and the elements of the group now act on the
product rule. % claim-ok: motif of this section, stated in the section plan; no external fact
The group ring $RG$ of a group $G$ over a commutative ring $R$ with
$1$ is the free $R$-module on $G$, with the product extended from $G$
\cite[p.~2956]{Spiegel1994}. Write $D$ for a derivation of $RG$. On a group
element the Leibniz rule reads $D(xy)=D(x)y+xD(y)$, so the values of $D$ on $G$
decide its values on whole of $RG$ by Lemma~\ref{lem:group-hopf:spiegel-extension}.

\begin{lemma}[{\cite[Lemma~1, p.~2956]{Spiegel1994}}]\label{lem:group-hopf:spiegel-extension}
Let $R$ be a commutative ring with $1$, $G$ a group, and $D\colon G\to RG$ a map with
$D(xy)=D(x)y+xD(y)$ for all $x,y\in G$. Then $D$ extends to a derivation of $RG$.
\end{lemma}

Spiegel proves the extension formula by linearity: write an element of $RG$ as a
finite sum $\sum r_g g$ and send it to $\sum r_g D(g)$ \cite[p.~2956]{Spiegel1994}.
He then asks which derivations of the integral group ring $\mathbb{Z}G$ occur, first
for a finite group (Theorem~\ref{thm:group-hopf:spiegel}) \cite[p.~2955]{Spiegel1994}.

\begin{theorem}[{\cite[Theorem~1, p.~2956]{Spiegel1994}}]\label{thm:group-hopf:spiegel}
Let $\mathbb{Z}G$ be the integral group ring of the finite group $G$. Every
derivation of $\mathbb{Z}G$ is inner.
\end{theorem}

The proof extends a derivation $D$ of $\mathbb{Z}G$ to the rational group algebra
$\mathbb{Q}G$ by Lemma~\ref{lem:group-hopf:spiegel-extension}. The theorem of
Maschke makes $\mathbb{Q}G$ semisimple \cite[p.~2956]{Spiegel1994}. On each simple component
the derivation is linear over the centre, so a consequence of the Noether--Skolem
theorem makes it inner \cite[pp.~2956--2957]{Spiegel1994}. Spiegel then clears
the fractions \cite[pp.~2957--2958]{Spiegel1994}: $D= D_f$ with $f\in\mathbb{Q}G$
and $D(h)\in\mathbb{Z}G$ for $h\in G$ give $(r_g-r_{h^{-1}gh})\in\mathbb{Z}$.
The fractional parts of the coefficients of $f$ form a central element, so $f$
may be taken inside $\mathbb{Z}G$ \cite[pp.~2957--2958]{Spiegel1994}. For an
abelian group $G$ there is more to say: $L(x)=D(x)x^{-1}$ is a homomorphism from
$G$ to the additive group of $\mathbb{Z}G$ \cite[p.~2958]{Spiegel1994}. The
image $L(G)$ is free abelian \cite[p.~2958]{Spiegel1994}. Take a basis of $L(G)$
and choose one element of $G$ above each basis element; these elements and
$\ker L$ give the direct product \cite[p.~2958]{Spiegel1994}.

\begin{theorem}[{\cite[Theorem~2, pp.~2958--2959]{Spiegel1994}}]\label{thm:group-hopf:spiegel-abelian}
Let $G$ be an abelian group.
\begin{enumerate}
\item If $D$ is a derivation of $\mathbb{Z}G$, there exist subgroups $A,B$ of $G$,
with $A$ a free abelian group having a basis $\{a_i\mid i\in I\}$, and a linearly
independent set $\{y_i\mid i\in I\}$ of the additive group of $\mathbb{Z}G$ such
that $G$ is the direct product of $A$ and $B$, and if $g\in G$ with $g=ab$, $a\in A$,
$b\in B$ and $a=a_{i_1}^{e_1}\cdots a_{i_n}^{e_n}$, then
$D(g)=(e_1 y_{i_1}+\cdots+e_n y_{i_n})g$.
\item Conversely, if $G$ is the direct product of groups $A,B$ with $A$ a free
abelian group with basis $\{a_i\mid i\in I\}$, and $\{y_i\mid i\in I\}$ is a set of
linearly independent elements of the additive group of $\mathbb{Z}G$, then the
mapping $D\colon G\to\mathbb{Z}G$ given by
$D(g)=(e_1 y_{i_1}+\cdots+e_n y_{i_n})g$ extends to a derivation of $\mathbb{Z}G$.
Here $g=ab$, $a\in A$, $b\in B$ and $a=a_{i_1}^{e_1}\cdots a_{i_n}^{e_n}$.
\end{enumerate}
\end{theorem}

So a derivation of $\mathbb{Z}G$ for $G$ abelian is a choice of a free factor of $G$
and of a free family in $\mathbb{Z}G$. Torsion-free factors give outer derivations
\cite[Theorem~2, pp.~2958--2959]{Spiegel1994}.

So the story now moves from $\mathbb{Z}$ to a semiprime coefficient ring. Ferrero,
Giambruno and Polcino Milies take a ring with unity $R$ and call a derivation
$d\colon RG\to RG$ an $R$-derivation when $d(R)=0$ \cite[p.~434]{FerreroGiambrunoPolcinoMilies1995}.
They also record a decomposition \cite[p.~434]{FerreroGiambrunoPolcinoMilies1995}:
a derivation $d$ of $RG$ with $d(R)\subseteq R$ writes as $d=d_1+d_2$. Here $d_1$
is a derivation of $R$ extended by $d_1(G)=0$, and $d_2$ is an $R$-derivation of
$RG$. Their main theorem puts a finiteness condition on the centre (Theorem~\ref{thm:group-hopf:fgp}).

\begin{theorem}[{\cite[Theorem~1.1, p.~434]{FerreroGiambrunoPolcinoMilies1995}}]\label{thm:group-hopf:fgp}
Let $R$ be a semiprime ring and $G$ a torsion group such that $[G:Z(G)]<\infty$,
where $Z(G)$ denotes the center of $G$. Suppose that either $\operatorname{char}R=0$
or for every characteristic $p$ of $R$, $p$ does not divide $o(g)$, for all $g\in G$.
Then every $R$-derivation of $RG$ is inner.
\end{theorem}

The proof first treats a finite group $G$
\cite[pp.~434--436]{FerreroGiambrunoPolcinoMilies1995}. The ring $R$ embeds into
a direct sum of algebraically closed fields, where the theorem of Maschke applies
\cite[p.~435]{FerreroGiambrunoPolcinoMilies1995}. The general case reduces to a
finite subgroup generated by the supports of the values of $d$
\cite[pp.~434--436]{FerreroGiambrunoPolcinoMilies1995}. Section~3 of the paper
shows the hypothesis on $G$ is needed.

\begin{example}[{\cite[\S3, p.~437]{FerreroGiambrunoPolcinoMilies1995}}]\label{ex:group-hopf:outer}
Let $H$ be a finite non-abelian $p$-group and let $G=\bigoplus_i H_i$ be the direct
sum of infinitely many copies of $H$. Then $G$ is a locally finite nilpotent group.
Let $F$ be any field, and choose a non-central element $x_i$ of $H_i$ for each $i$.
The map $d\colon FG\to FG$ with $d(a)=\sum_{i\ge1}[x_i,a]$ is a derivation of $FG$:
on $FG_n$ it is the inner derivation induced by $x_1+\cdots+x_n$. It is not an inner
derivation of $FG$, since $d(H_i)\neq0$ for all indices $i$.
\end{example}

So an infinite group can carry a derivation that is inner on every finitely
generated subgroup but outer on the whole group.

Manju and Sharma count the inner derivations over any field, and they read the
count from the class sums \cite[p.~5]{ManjuSharma2026}. The class sums
form a basis of the centre of $FG$ \cite[p.~5]{ManjuSharma2026}. The proof
extends that basis to a basis of $FG$
\cite[p.~5]{ManjuSharma2026}. The inner derivations indexed by the added
elements then form a basis of the inner ones \cite[p.~5]{ManjuSharma2026}.

\begin{theorem}[{\cite[Theorem~2.15, p.~5]{ManjuSharma2026}}]\label{thm:group-hopf:manju-sharma-dim}
Let $F$ be an arbitrary field of any characteristic and $G$ a finite group of
order $n$ having $r$ conjugacy classes. Then the dimension of the space of inner
derivations of $FG$ is $n-r$. If $C_1,\dots,C_r$ are the distinct conjugacy
classes of $G$ with representatives $x_1,\dots,x_r$ respectively and $s$ is a
positive integer such that $|C_i|=1$ for all $i\in\{1,\dots,s\}$ and $|C_i|\ge2$
for all $i\in\{s+1,\dots,r\}$, then the set
$B_0=\{d_g\mid g\in\bigcup_{i=s+1}^{r}(C_i\setminus\{x_i\})\}$ is an $F$-basis of
the space of inner derivations of $FG$.
\end{theorem}

Here $d_g$ is the inner derivation $d_g(\alpha)=\alpha g-g\alpha$
\cite[Definition~2.2, p.~3]{ManjuSharma2026}. So the count of
Theorem~\ref{thm:group-hopf:manju-sharma-dim} is fixed by the class number
alone, for every finite group \cite[p.~5]{ManjuSharma2026}. The paper then works
the count out
for three families and asks when every derivation is inner
\cite[Corollary~3.6, p.~10]{ManjuSharma2026}. The dihedral group is
$D_{2n}=\langle a,b\mid a^n=1,\ b^2=1,\ (ab)^2=1\rangle$ of order $2n$, with
$n\ge3$ (Corollary~\ref{cor:group-hopf:manju-sharma-dihedral}) \cite[p.~6]{ManjuSharma2026}.

\begin{corollary}[{\cite[Corollary~3.6, p.~10]{ManjuSharma2026}}]\label{cor:group-hopf:manju-sharma-dihedral}
Let $F$ be an algebraic extension of a prime field and $p$ an odd rational
prime.
\begin{enumerate}
\item If $\operatorname{char}F$ is either $0$ or $p$ with $\gcd(n,p)=1$, then
every derivation of $FD_{2n}$ is inner.
\item If $\operatorname{char}F=p$ with $\gcd(n,p)\neq1$, then the inner
derivations form a proper subspace of the space of all derivations of $FD_{2n}$.
\end{enumerate}
\end{corollary}

The authors prove the two cases by counting both spaces
\cite[Theorem~3.5, p.~10]{ManjuSharma2026}. Their count of the inner derivations
of $FD_{2n}$ is $3\lfloor (n-1)/2\rfloor$ over any field
\cite[Theorem~3.5, p.~10]{ManjuSharma2026}. Their count of all derivations of
$FD_{2n}$ equals it exactly when $\operatorname{char}F$ is $0$ or a prime not
dividing $n$ \cite[Theorem~3.3, p.~7]{ManjuSharma2026}. A prime dividing the
group order therefore leaves room for outer derivations on a finite group
algebra \cite[Corollary~3.6, p.~10]{ManjuSharma2026}.

The derivation algebra itself carries a grading. Arutyunov and Zhiltsov
identify a derivation with a character of the groupoid of the adjoint
action \cite[p.~1]{ArutyunovZhiltsov2023}. Fix a normal subgroup $N$ of $G$
with $G/N$ abelian. Then the derivations split into pieces indexed by the
cosets of $N$ (Theorem~\ref{thm:group-hopf:grading}) \cite[p.~1]{ArutyunovZhiltsov2023}.

\begin{theorem}[{\cite[Theorem~1, pp.~1, 8]{ArutyunovZhiltsov2023}}]\label{thm:group-hopf:grading}
Let $G$ be a group and let $N$ be a normal subgroup of $G$ with $G/N$
abelian and $|G/N|>1$. Then the derivation algebra of the group algebra is
graded by $G/N$:
\[
\operatorname{Der}=\bigoplus_{k\in G/N}\operatorname{Der}_{k},
\qquad
[\operatorname{Der}_{k},\operatorname{Der}_{l}]\subset\operatorname{Der}_{kl}
\quad\text{for all }k,l\in G/N .
\]
\end{theorem}

The piece $\operatorname{Der}_{k}$ collects the derivations whose character
is supported on the coset $k$ \cite[p.~1]{ArutyunovZhiltsov2023}. Every
group that is not perfect has a nontrivial grading, because such a group has
a proper normal subgroup with abelian quotient
\cite[p.~1]{ArutyunovZhiltsov2023}.

Arutyunov and Naianzin carry the derivation problem to a family of groups
that they call BC-groups \cite[p.~2]{ArutyunovNaianzin2023}. The defining
condition bounds the diameter of the conjugate of any finite set
\cite[p.~2]{ArutyunovNaianzin2023}. For such a group every bounded
derivation from the group ring into an $\ell^{q}$ space is inner
\cite[Theorem~2, p.~2]{ArutyunovNaianzin2023}. This extends the case of
Johnson, who settled the derivations of $\ell^{1}(G)$ for a discrete group
\cite[p.~2]{ArutyunovNaianzin2023}.

The question has a block form for finite groups. Let $kG$ be a group algebra
over a field $k$ of characteristic $p$, and let $B$ be a block of $kG$ with
defect group $P$ \cite[p.~2]{AntrobusBriggsRubioDegrassi2026}. Linckelmann
made the conjecture that the first Hochschild cohomology of $B$ is nonzero
whenever $P$ is nontrivial \cite[Conjecture, p.~2]{AntrobusBriggsRubioDegrassi2026}.
Antrobus, Briggs and Rubio y Degrassi prove many cases of it (Theorem~\ref{thm:group-hopf:block-dimension})
\cite[Theorem~B, p.~2]{AntrobusBriggsRubioDegrassi2026}.\begin{theorem}[{\cite[Theorem~A, p.~2]{AntrobusBriggsRubioDegrassi2026}}]\label{thm:group-hopf:block-dimension}
Let $B$ be a block of $kG$ with defect group $P$. Then the dimension of
$HH^{1}(B)$ is at least the number of conjugacy classes of extra-strong
non-Schur elements of $G$ relative to $P$.
\end{theorem}

The cases include the principal blocks with abelian defect group, and all
blocks of the symmetric and alternating groups. They include the blocks of
finite groups of Lie type in defining characteristic, and the blocks of
general linear groups in any characteristic
\cite[Theorem~B, p.~2]{AntrobusBriggsRubioDegrassi2026}.
When $p>5$ the conjecture holds for all blocks of maximal defect
\cite[Theorem~C, p.~3]{AntrobusBriggsRubioDegrassi2026}.

\section{Harmonic analysis: derivations on group rings}\label{sec:group-hopf:harmonic-analysis}

Fox built the harmonic analysis of a group ring in 1953 \cite{Fox1953}. The
classical theory begins with a definition of derivation that is the Leibniz
rule on the group itself
\cite[equations (1.1)--(1.2), p.~549]{Fox1953}.

\begin{definition}[{\cite[equations (1.1)--(1.2), p.~549]{Fox1953}}]\label{def:group-hopf:fox-derivation}
Let $JG$ be the group ring of a group $G$ over a commutative ring $J$ with
unit. By a \emph{derivation} in $JG$ is meant any mapping $D$ of $JG$ into
itself which satisfies
\[
D(u+v)=Du+Dv,\qquad D(u\ast v)=Du\cdot v^{0}+u\ast Dv\quad\text{for }u,v\in JG,
\]
where $v^{0}\in J$ is the coefficient sum of $v$.
\end{definition}

For elements of the group itself, the second law takes the simpler form
$D(gh)=Dg+g\,Dh$ \cite[(1.2)', p.~549]{Fox1953}, which is the Leibniz rule for
the group multiplication. Fox shows the set of derivations in $JG$ forms a
right $JG$-module, with $(D\cdot v)(u)=Du\,v$ \cite[p.~549]{Fox1953}. On the
free group ring, the structure is fully described by (Theorem~\ref{thm:group-hopf:fox-free-derivatives}) \cite[Theorem, equations (2.1) and (2.2), p.~550]{Fox1953}.

\begin{theorem}[{\cite[Theorem, equations (2.1) and (2.2), p.~550]{Fox1953}}]\label{thm:group-hopf:fox-free-derivatives}
Let $X$ be a free group with generators $(x)=(x_1,x_2,\dots)$ and let $JX$ be
its group ring. To each generator $x_j$ of $X$ there corresponds a derivation
$f(x)\mapsto D_j f(x)=f_{x_j}(x)=\partial f(x)/\partial x_j$, called the
\emph{derivative with respect to} $x_j$ \cite[p.~550]{Fox1953}, which has the property
$\partial x_k/\partial x_j=\delta_{jk}$. Furthermore, there is one and only one
derivation $f(x)\mapsto f'(x)$ mapping $x_1,x_2,\ldots$ into prescribed
elements $h_1(x),h_2(x),\ldots$ of $JX$; it is given by
$f'(x)=\sum_j f_{x_j}(x)\cdot h_j(x)$.
\end{theorem}

On the continuous side, the same rule governs the operator algebras of
Chapter~\ref{ch:banach-operator}, where each derivation of a von Neumann
algebra is inner \cite[Theorem~15, p.~290]{Kadison1966}.

Section~\ref{sec:group-hopf:hopf} moves from the group ring to a Hopf algebra,
where the coproduct itself asks a derivation to respect a second operation.

A continuous linear map of the group algebra $L^1(G)$ into the measure algebra is determined by its action on one element, which is the determination question for the group-algebra derivations of the chapter~\cite{WlLinearMapsOnGroupAlgebrasDeter2019}.

For a finite cyclic group $G$ and a pair $(\sigma,\tau)$ of endomorphisms, the inner and outer twisted derivations of the group ring $R[G]$ are determined~\cite{WlInnerAndOuterTwistedDerivation2019}.
